% ===== PRÉAMBULE CANONIQUE — à coller VERBATIM en tête de CHAQUE .tex =====
\documentclass[11pt,a4paper]{article}
\usepackage[utf8]{inputenc}\usepackage[T1]{fontenc}\usepackage{lmodern}
\usepackage[french]{babel}
\usepackage{amsmath,amssymb,mathtools}\usepackage{icomma}
\usepackage{siunitx}\sisetup{locale=FR}  % \qty{}{}, \si{}, \ang{}
\usepackage{xcolor}\usepackage{colortbl}
\usepackage{tikz}\usepackage{pgfplots}\pgfplotsset{compat=1.18}
\usepackage[siunitx,europeanresistors]{circuitikz} % circuits (élec.)
\usepackage[most]{tcolorbox}\usepackage{enumitem}\usepackage{array,booktabs}
\usepackage[a4paper,margin=2.2cm]{geometry}
\usetikzlibrary{arrows.meta,calc,angles,quotes,positioning,patterns,%
  backgrounds,shapes.geometric,decorations.pathmorphing,decorations.markings,intersections}
\definecolor{primary}{RGB}{222,49,99}\definecolor{primarydark}{RGB}{150,22,55}
\definecolor{defcol}{RGB}{0,114,178}\definecolor{methcol}{RGB}{52,92,148}
\definecolor{excol}{RGB}{0,135,90}\definecolor{attncol}{RGB}{200,40,55}
\tcbset{boxbase/.style={enhanced,breakable,boxrule=0.9pt,arc=2.5pt,
  left=3mm,right=3mm,top=2.3mm,bottom=2.3mm,fonttitle=\bfseries,coltitle=white}}
% --- boîtes TUTORIEL (noms EXACTS, ne pas renommer) ---
\newtcolorbox{cle}[1][]{boxbase,colback=defcol!6,colframe=defcol,
  title={\textsc{À retenir}\ifx\relax#1\relax\relax\else\textmd{ : \itshape #1}\fi}}
\newtcolorbox{methode}[1][]{boxbase,colback=methcol!7,colframe=methcol,sharp corners,
  title={\textsc{Méthode}\ifx\relax#1\relax\relax\else\textmd{ --- \itshape #1}\fi}}
\newtcolorbox{exemplebox}[1][]{boxbase,colback=excol!5,colframe=excol,
  title={\textsc{Exemple}\ifx\relax#1\relax\relax\else\textmd{ : \itshape #1}\fi}}
\newtcolorbox{piege}[1][]{boxbase,colback=attncol!6,colframe=attncol,
  title={\textsc{Piège}\ifx\relax#1\relax\relax\else\textmd{ : \itshape #1}\fi}}
\newtcolorbox{remarque}[1][]{boxbase,colback=black!4,colframe=black!45,
  title={\textsc{Remarque}\ifx\relax#1\relax\relax\else\textmd{ : \itshape #1}\fi}}
% --- boîtes EXERCICE / CORRIGÉ (noms EXACTS, auto-comptées) ---
\newtcolorbox[auto counter]{exo}[1][]{boxbase,colback=primary!4,colframe=primary,
  title={\textsc{Exercice}~\thetcbcounter\ifx\relax#1\relax\relax\else\textmd{ --- \itshape #1}\fi}}
\newtcolorbox[auto counter]{corrige}[1][]{boxbase,colback=excol!4,colframe=excol,
  title={\textsc{Corrigé}~\thetcbcounter\ifx\relax#1\relax\relax\else\textmd{ --- \itshape #1}\fi}}
\newcommand{\vv}[1]{\vec{#1}}\newcommand{\ds}{\displaystyle}
\newcommand{\R}{\mathbb{R}}\newcommand{\N}{\mathbb{N}}\newcommand{\Z}{\mathbb{Z}}
% =========================================================================
\begin{document}
\sisetup{per-mode=symbol}

\begin{corrige}[chute non libre : l'accélération]
\begin{enumerate}[label=\alph*)]
  \item Dans le sens de la chute : le poids $G=mg$ (vers le bas) et la résistance $F_f$ (vers le haut),
        donc $G-F_f=ma$.
  \item $a=\dfrac{mg-F_f}{m}=g-\dfrac{F_f}{m}=10-\dfrac{300}{80}=10-\num{3,75}=\qty{6,25}{\metre\per\second\squared}$.
\end{enumerate}
\end{corrige}

\begin{corrige}[retrouver la résistance de l'air]
De $g-\dfrac{F_f}{m}=a$, on tire $F_f=m(g-a)=2\times(10-8)=\qty{4}{\newton}$.
\end{corrige}

\begin{corrige}[la pomme de pin à vitesse constante]
\begin{enumerate}[label=\alph*)]
  \item Vitesse constante $\Rightarrow a=0$. D'après la $1^{\text{re}}$ loi, les forces se compensent :
        $F_f=G$.
  \item $F_f=G=mg=\num{0,05}\times 10=\qty{0,5}{\newton}$ (avec $m=\qty{50}{\gram}=\qty{0,05}{\kilogram}$).
\end{enumerate}
\end{corrige}

\begin{corrige}[pourquoi $a<g$ ?]
\begin{enumerate}[label=\alph*)]
  \item $a=\dfrac{G-F_f}{m}=\dfrac{20-4}{2}=\qty{8}{\metre\per\second\squared}$.
  \item $a=\qty{8}{\metre\per\second\squared}<g=\qty{10}{\metre\per\second\squared}$ : la résistance de l'air
        retranche $\tfrac{F_f}{m}=\qty{2}{\metre\per\second\squared}$. Sans air ($F_f=0$) on retrouverait $a=g$.
\end{enumerate}
\end{corrige}

\begin{corrige}[deux masses lancées sur la Lune]
\begin{enumerate}[label=\alph*)]
  \item Oui. Sans air, chaque masse n'est soumise qu'à son poids : $a=\dfrac{mg}{m}=g=\qty{1,6}{\metre\per\second\squared}$,
        \emph{indépendante de la masse}.
  \item Elles montent \textbf{à la même hauteur} : même accélération et même vitesse initiale $\Rightarrow$
        mouvement identique. La masse n'intervient pas.
\end{enumerate}
\end{corrige}
\end{document}
